\section{Semantics}

% BEGIN GENERATED ARTIFACT: def:truth-assignment-propositional-logic
\begin{tcolorbox}[colback=propbox, colframe=propborder, arc=2pt,
  left=6pt, right=6pt, top=4pt, bottom=4pt,
  title={\small\textbf{Definition (Truth Assignment)}},
  fonttitle=\small\bfseries]
\begin{definition}[Truth Assignment]
\label{def:truth-assignment-propositional-logic}
A truth assignment for propositional logic is a function
\[
v:\Prop\to\mathbb{B}.
\]
Thus a truth assignment assigns either \(\mathrm{False}\) or \(\mathrm{True}\) to each propositional variable.
\end{definition}
\end{tcolorbox}

\begin{remark*}[Standard quantified statement]
\[
\operatorname{TruthAssignment}(v)
\Longleftrightarrow
v:\Prop\to\mathbb{B}.
\]
\end{remark*}

\begin{remark*}[Definition predicate reading]
\[
\operatorname{TruthAssignment}(v)
\]
means that \(v\) assigns a truth value to every propositional variable.
\end{remark*}

\begin{remark*}[Interpretation]
A truth assignment fixes the truth values of atomic formulas. The recursive semantics of the connectives then determines the truth value of each molecular formula.
\end{remark*}

\begin{dependencies}
\begin{itemize}
  \item \hyperref[def:truth-value-propositional-logic]{Truth Value}
  \item \hyperref[def:propositional-variable]{Propositional Variable}
  \item \hyperref[def:well-formed-formula]{Well-Formed Formula}
\end{itemize}
\end{dependencies}
% END GENERATED ARTIFACT: def:truth-assignment-propositional-logic
% BEGIN GENERATED ARTIFACT: thm:unique-extension-truth-assignment-propositional-logic
\begin{tcolorbox}[colback=thmbox, colframe=thmborder, arc=2pt,
  left=6pt, right=6pt, top=4pt, bottom=4pt,
  title={\small\textbf{Theorem (Unique Extension of a Truth Assignment)}},
  fonttitle=\small\bfseries]
\begin{theorem}[Unique Extension of a Truth Assignment]
\label{thm:unique-extension-truth-assignment-propositional-logic}
Every truth assignment \(v:\Prop\to\mathbb{B}\) extends uniquely to a function
\[
\widehat v:\WFF\to\mathbb{B}
\]
satisfying the following clauses:
\[
\widehat v(P)=v(P)\quad(P\in\Prop),
\]
\[
\widehat v(\neg\varphi)=\mathrm{True}
\Longleftrightarrow
\widehat v(\varphi)=\mathrm{False},
\]
\[
\widehat v(\varphi\land\psi)=\mathrm{True}
\Longleftrightarrow
\widehat v(\varphi)=\mathrm{True}
\text{ and }
\widehat v(\psi)=\mathrm{True},
\]
\[
\widehat v(\varphi\lor\psi)=\mathrm{True}
\Longleftrightarrow
\widehat v(\varphi)=\mathrm{True}
\text{ or }
\widehat v(\psi)=\mathrm{True},
\]
\[
\widehat v(\varphi\to\psi)=\mathrm{True}
\Longleftrightarrow
\widehat v(\varphi)=\mathrm{False}
\text{ or }
\widehat v(\psi)=\mathrm{True},
\]
and
\[
\widehat v(\varphi\leftrightarrow\psi)=\mathrm{True}
\Longleftrightarrow
\widehat v(\varphi)=\widehat v(\psi).
\]

\noindent\hyperref[prf:unique-extension-truth-assignment-propositional-logic]{Go to proof.}
\end{theorem}
\end{tcolorbox}

\begin{remark*}[Standard quantified statement]
\[
\begin{aligned}
&(\forall v:\Prop\to\mathbb{B})(\exists!\widehat v:\WFF\to\mathbb{B})\\
&\quad\text{such that, for all }P\in\Prop\text{ and }\varphi,\psi\in\WFF,\\
&\quad \widehat v(P)=v(P),\\
&\quad \widehat v(\neg\varphi)=\mathrm{True}
   \Longleftrightarrow \widehat v(\varphi)=\mathrm{False},\\
&\quad \widehat v(\varphi\land\psi)=\mathrm{True}
   \Longleftrightarrow
   (\widehat v(\varphi)=\mathrm{True}\land \widehat v(\psi)=\mathrm{True}),\\
&\quad \widehat v(\varphi\lor\psi)=\mathrm{True}
   \Longleftrightarrow
   (\widehat v(\varphi)=\mathrm{True}\lor \widehat v(\psi)=\mathrm{True}),\\
&\quad \widehat v(\varphi\to\psi)=\mathrm{True}
   \Longleftrightarrow
   (\widehat v(\varphi)=\mathrm{False}\lor \widehat v(\psi)=\mathrm{True}),\\
&\quad \widehat v(\varphi\leftrightarrow\psi)=\mathrm{True}
   \Longleftrightarrow
   \widehat v(\varphi)=\widehat v(\psi).
\end{aligned}
\]
\end{remark*}

\begin{remark*}[Interpretation]
A valuation on variables determines exactly one truth value for every formula. This theorem is the semantic counterpart of structural recursion and unique decomposition.
\end{remark*}

\begin{dependencies}
\begin{itemize}
  \item \hyperref[def:truth-assignment-propositional-logic]{Truth Assignment}
  \item \hyperref[def:well-formed-formula]{Well-Formed Formula}
  \item \hyperref[thm:structural-recursion-propositional-formulas]{Structural Recursion for Propositional Formulas}
  \item \hyperref[thm:unique-decomposition-propositional-formulas]{Unique Decomposition for Propositional Formulas}
\end{itemize}
\end{dependencies}
% END GENERATED ARTIFACT: thm:unique-extension-truth-assignment-propositional-logic
% BEGIN GENERATED ARTIFACT: def:satisfaction-propositional-formula
\begin{tcolorbox}[colback=propbox, colframe=propborder, arc=2pt,
  left=6pt, right=6pt, top=4pt, bottom=4pt,
  title={\small\textbf{Definition (Satisfaction of a Formula)}},
  fonttitle=\small\bfseries]
\begin{definition}[Satisfaction of a Formula]
\label{def:satisfaction-propositional-formula}
Let \(v\) be a truth assignment and let \(\varphi\in\WFF\). We say that \(v\) satisfies \(\varphi\), written
\[
v\models\varphi,
\]
if
\[
\widehat v(\varphi)=\mathrm{True}.
\]
\end{definition}
\end{tcolorbox}

\begin{remark*}[Standard quantified statement]
\[
v\models\varphi
\Longleftrightarrow
\operatorname{TruthAssignment}(v)
\land
\varphi\in\WFF
\land
\widehat v(\varphi)=\mathrm{True}.
\]
\end{remark*}

\begin{remark*}[Definition predicate reading]
\[
\operatorname{SatisfiesFormula}(v,\varphi)
\]
means \(v\models\varphi\).
\end{remark*}

\begin{remark*}[Negated quantified statement]
\[
v\not\models\varphi
\Longleftrightarrow
\widehat v(\varphi)=\mathrm{False}.
\]
\end{remark*}

\begin{remark*}[Interpretation]
The notation \(v\models\varphi\) means that the formula \(\varphi\) is true under the truth assignment \(v\).
\end{remark*}

\begin{dependencies}
\begin{itemize}
  \item \hyperref[def:truth-assignment-propositional-logic]{Truth Assignment}
  \item \hyperref[thm:unique-extension-truth-assignment-propositional-logic]{Unique Extension of a Truth Assignment}
  \item \hyperref[def:well-formed-formula]{Well-Formed Formula}
\end{itemize}
\end{dependencies}
% END GENERATED ARTIFACT: def:satisfaction-propositional-formula
% BEGIN GENERATED ARTIFACT: def:satisfaction-set-propositional-logic
\begin{tcolorbox}[colback=propbox, colframe=propborder, arc=2pt,
  left=6pt, right=6pt, top=4pt, bottom=4pt,
  title={\small\textbf{Definition (Satisfaction of a Set of Formulas)}},
  fonttitle=\small\bfseries]
\begin{definition}[Satisfaction of a Set of Formulas]
\label{def:satisfaction-set-propositional-logic}
Let \(\Gamma\subseteq\WFF\). A truth assignment \(v\) satisfies \(\Gamma\), written
\[
v\models\Gamma,
\]
if \(v\models\gamma\) for every \(\gamma\in\Gamma\).
\end{definition}
\end{tcolorbox}

\begin{remark*}[Standard quantified statement]
\[
v\models\Gamma
\Longleftrightarrow
\operatorname{TruthAssignment}(v)
\land
\Gamma\subseteq\WFF
\land
(\forall\gamma\in\Gamma)(v\models\gamma).
\]
\end{remark*}

\begin{remark*}[Definition predicate reading]
\[
\operatorname{SatisfiesSet}(v,\Gamma)
\]
means \(v\) satisfies every formula in \(\Gamma\).
\end{remark*}

\begin{remark*}[Negated quantified statement]
\[
v\not\models\Gamma
\Longleftrightarrow
\exists\gamma\in\Gamma\text{ such that }v\not\models\gamma.
\]
\end{remark*}

\begin{remark*}[Interpretation]
Satisfaction of a set means simultaneous satisfaction of every formula in the set.
\end{remark*}

\begin{dependencies}
\begin{itemize}
  \item \hyperref[def:satisfaction-propositional-formula]{Satisfaction of a Formula}
  \item \hyperref[def:truth-assignment-propositional-logic]{Truth Assignment}
\end{itemize}
\end{dependencies}
% END GENERATED ARTIFACT: def:satisfaction-set-propositional-logic
% BEGIN GENERATED ARTIFACT: def:truth-table-propositional-logic
\begin{tcolorbox}[colback=propbox, colframe=propborder, arc=2pt,
  left=6pt, right=6pt, top=4pt, bottom=4pt,
  title={\small\textbf{Definition (Truth Table)}},
  fonttitle=\small\bfseries]
\begin{definition}[Truth Table]
\label{def:truth-table-propositional-logic}
Let \(\varphi\in\WFF\) and suppose the variables occurring in \(\varphi\) are among \(P_1,\ldots,P_n\). A truth table for \(\varphi\) is a finite table with one row for each assignment
\[
a:\{P_1,\ldots,P_n\}\to\mathbb{B}
\]
and a final column recording the value of \(\widehat v(\varphi)\) for any truth assignment \(v\) extending \(a\).
\end{definition}
\end{tcolorbox}

\begin{remark*}[Standard quantified statement]
\[
\operatorname{TruthTable}(\varphi)
\Longleftrightarrow
\varphi\in\WFF
\land
\operatorname{Var}(\varphi)=\{P_1,\ldots,P_n\}
\land
\text{the table has }2^n\text{ assignment rows}.
\]
\end{remark*}

\begin{remark*}[Interpretation]
A truth table is finite because every formula contains only finitely many variables. It records the complete truth behavior of a formula.
\end{remark*}

\begin{dependencies}
\begin{itemize}
  \item \hyperref[def:formula-variable-set-propositional-logic]{Variable Set of a Formula}
  \item \hyperref[lem:finiteness-of-variables-propositional-formulas]{Finiteness of Variables in a Formula}
  \item \hyperref[thm:unique-extension-truth-assignment-propositional-logic]{Unique Extension of a Truth Assignment}
\end{itemize}
\end{dependencies}
% END GENERATED ARTIFACT: def:truth-table-propositional-logic
% BEGIN GENERATED ARTIFACT: def:tautology-propositional-logic
\begin{tcolorbox}[colback=propbox, colframe=propborder, arc=2pt,
  left=6pt, right=6pt, top=4pt, bottom=4pt,
  title={\small\textbf{Definition (Tautology)}},
  fonttitle=\small\bfseries]
\begin{definition}[Tautology]
\label{def:tautology-propositional-logic}
A formula \(\varphi\in\WFF\) is a tautology, written
\[
\models\varphi,
\]
if every truth assignment satisfies \(\varphi\).
\end{definition}
\end{tcolorbox}

\begin{remark*}[Standard quantified statement]
\[
\models\varphi
\Longleftrightarrow
\varphi\in\WFF
\land
(\forall v:\Prop\to\mathbb{B})(v\models\varphi).
\]
\end{remark*}

\begin{remark*}[Definition predicate reading]
\[
\operatorname{Tautology}(\varphi)
\]
means \(\models\varphi\).
\end{remark*}

\begin{remark*}[Interpretation]
A tautology is true under every possible assignment of truth values to propositional variables.
\end{remark*}

\begin{dependencies}
\begin{itemize}
  \item \hyperref[def:satisfaction-propositional-formula]{Satisfaction of a Formula}
  \item \hyperref[def:truth-assignment-propositional-logic]{Truth Assignment}
\end{itemize}
\end{dependencies}
% END GENERATED ARTIFACT: def:tautology-propositional-logic
% BEGIN GENERATED ARTIFACT: def:contradiction-propositional-logic
\begin{tcolorbox}[colback=propbox, colframe=propborder, arc=2pt,
  left=6pt, right=6pt, top=4pt, bottom=4pt,
  title={\small\textbf{Definition (Contradiction)}},
  fonttitle=\small\bfseries]
\begin{definition}[Contradiction]
\label{def:contradiction-propositional-logic}
A formula \(\varphi\in\WFF\) is a contradiction if no truth assignment satisfies \(\varphi\).
\end{definition}
\end{tcolorbox}

\begin{remark*}[Standard quantified statement]
\[
\operatorname{Contradiction}(\varphi)
\Longleftrightarrow
\varphi\in\WFF
\land
(\forall v:\Prop\to\mathbb{B})(v\not\models\varphi).
\]
\end{remark*}

\begin{remark*}[Definition predicate reading]
\[
\operatorname{Contradiction}(\varphi)
\]
means that \(\varphi\) is false under every truth assignment.
\end{remark*}

\begin{remark*}[Interpretation]
A contradiction has no true row in its truth table.
\end{remark*}

\begin{dependencies}
\begin{itemize}
  \item \hyperref[def:satisfaction-propositional-formula]{Satisfaction of a Formula}
  \item \hyperref[def:truth-assignment-propositional-logic]{Truth Assignment}
\end{itemize}
\end{dependencies}
% END GENERATED ARTIFACT: def:contradiction-propositional-logic
% BEGIN GENERATED ARTIFACT: def:satisfiable-formula-propositional-logic
\begin{tcolorbox}[colback=propbox, colframe=propborder, arc=2pt,
  left=6pt, right=6pt, top=4pt, bottom=4pt,
  title={\small\textbf{Definition (Satisfiable Formula)}},
  fonttitle=\small\bfseries]
\begin{definition}[Satisfiable Formula]
\label{def:satisfiable-formula-propositional-logic}
A formula \(\varphi\in\WFF\) is satisfiable if at least one truth assignment satisfies it.
\end{definition}
\end{tcolorbox}

\begin{remark*}[Standard quantified statement]
\[
\operatorname{SatisfiableFormula}(\varphi)
\Longleftrightarrow
\varphi\in\WFF
\land
\exists v:\Prop\to\mathbb{B}\text{ such that }v\models\varphi.
\]
\end{remark*}

\begin{remark*}[Definition predicate reading]
\[
\operatorname{SatisfiableFormula}(\varphi)
\]
means that \(\varphi\) is true under at least one truth assignment.
\end{remark*}

\begin{remark*}[Interpretation]
A satisfiable formula has at least one true row in its truth table.
\end{remark*}

\begin{dependencies}
\begin{itemize}
  \item \hyperref[def:satisfaction-propositional-formula]{Satisfaction of a Formula}
  \item \hyperref[def:truth-assignment-propositional-logic]{Truth Assignment}
\end{itemize}
\end{dependencies}
% END GENERATED ARTIFACT: def:satisfiable-formula-propositional-logic
% BEGIN GENERATED ARTIFACT: def:contingency-propositional-logic
\begin{tcolorbox}[colback=propbox, colframe=propborder, arc=2pt,
  left=6pt, right=6pt, top=4pt, bottom=4pt,
  title={\small\textbf{Definition (Contingency)}},
  fonttitle=\small\bfseries]
\begin{definition}[Contingency]
\label{def:contingency-propositional-logic}
A formula \(\varphi\in\WFF\) is a contingency if it is satisfiable but not a tautology.
\end{definition}
\end{tcolorbox}

\begin{remark*}[Standard quantified statement]
\[
\operatorname{Contingency}(\varphi)
\Longleftrightarrow
\operatorname{SatisfiableFormula}(\varphi)
\land
\neg \operatorname{Tautology}(\varphi).
\]
\end{remark*}

\begin{remark*}[Definition predicate reading]
\[
\operatorname{Contingency}(\varphi)
\]
means that \(\varphi\) is true under some truth assignments and false under others.
\end{remark*}

\begin{remark*}[Interpretation]
A contingency has at least one true row and at least one false row in its truth table.
\end{remark*}

\begin{dependencies}
\begin{itemize}
  \item \hyperref[def:satisfiable-formula-propositional-logic]{Satisfiable Formula}
  \item \hyperref[def:tautology-propositional-logic]{Tautology}
\end{itemize}
\end{dependencies}
% END GENERATED ARTIFACT: def:contingency-propositional-logic
% BEGIN GENERATED ARTIFACT: def:logical-consequence-propositional-logic
\begin{tcolorbox}[colback=propbox, colframe=propborder, arc=2pt,
  left=6pt, right=6pt, top=4pt, bottom=4pt,
  title={\small\textbf{Definition (Logical Consequence)}},
  fonttitle=\small\bfseries]
\begin{definition}[Logical Consequence]
\label{def:logical-consequence-propositional-logic}
Let \(\Gamma\subseteq\WFF\) and \(\varphi\in\WFF\). The formula \(\varphi\) is a logical consequence of \(\Gamma\), written
\[
\Gamma\models\varphi,
\]
if every truth assignment satisfying \(\Gamma\) also satisfies \(\varphi\).
\end{definition}
\end{tcolorbox}

\begin{remark*}[Standard quantified statement]
\[
\Gamma\models\varphi
\Longleftrightarrow
\Gamma\subseteq\WFF
\land
\varphi\in\WFF
\land
(\forall v:\Prop\to\mathbb{B})
\left[
v\models\Gamma\to v\models\varphi
\right].
\]
\end{remark*}

\begin{remark*}[Definition predicate reading]
\[
\operatorname{LogicalConsequence}(\Gamma,\varphi)
\]
means \(\Gamma\models\varphi\).
\end{remark*}

\begin{remark*}[Negated quantified statement]
\[
\Gamma\not\models\varphi
\Longleftrightarrow
\exists v:\Prop\to\mathbb{B}
\left[
v\models\Gamma
\land
v\not\models\varphi
\right].
\]
\end{remark*}

\begin{remark*}[Interpretation]
Logical consequence means truth preservation from premises to conclusion. There is no truth assignment under which all premises are true and the conclusion is false.
\end{remark*}

\begin{dependencies}
\begin{itemize}
  \item \hyperref[def:satisfaction-set-propositional-logic]{Satisfaction of a Set of Formulas}
  \item \hyperref[def:satisfaction-propositional-formula]{Satisfaction of a Formula}
\end{itemize}
\end{dependencies}
% END GENERATED ARTIFACT: def:logical-consequence-propositional-logic
% BEGIN GENERATED ARTIFACT: def:logical-equivalence-propositional-logic
\begin{tcolorbox}[colback=propbox, colframe=propborder, arc=2pt,
  left=6pt, right=6pt, top=4pt, bottom=4pt,
  title={\small\textbf{Definition (Logical Equivalence)}},
  fonttitle=\small\bfseries]
\begin{definition}[Logical Equivalence]
\label{def:logical-equivalence-propositional-logic}
Formulas \(\varphi,\psi\in\WFF\) are logically equivalent, written
\[
\varphi\equiv\psi,
\]
if they have the same truth value under every truth assignment.
\end{definition}
\end{tcolorbox}

\begin{remark*}[Standard quantified statement]
\[
\varphi\equiv\psi
\Longleftrightarrow
\varphi,\psi\in\WFF
\land
(\forall v:\Prop\to\mathbb{B})
\left[
\widehat v(\varphi)=\widehat v(\psi)
\right].
\]
\end{remark*}

\begin{remark*}[Definition predicate reading]
\[
\operatorname{LogicalEquivalence}(\varphi,\psi)
\]
means \(\varphi\equiv\psi\).
\end{remark*}

\begin{remark*}[Interpretation]
Logically equivalent formulas have identical truth-table columns. They are interchangeable for semantic purposes.
\end{remark*}

\begin{dependencies}
\begin{itemize}
  \item \hyperref[thm:unique-extension-truth-assignment-propositional-logic]{Unique Extension of a Truth Assignment}
  \item \hyperref[def:truth-assignment-propositional-logic]{Truth Assignment}
\end{itemize}
\end{dependencies}
% END GENERATED ARTIFACT: def:logical-equivalence-propositional-logic
% BEGIN GENERATED ARTIFACT: lem:coincidence-lemma-propositional-logic
\begin{tcolorbox}[colback=lembox, colframe=lemborder, arc=2pt,
  left=6pt, right=6pt, top=4pt, bottom=4pt,
  title={\small\textbf{Lemma (Coincidence Lemma for Propositional Logic)}},
  fonttitle=\small\bfseries]
\begin{lemma}[Coincidence Lemma for Propositional Logic]
\label{lem:coincidence-lemma-propositional-logic}
Let \(v,w:\Prop\to\mathbb{B}\) be truth assignments. If \(v\) and \(w\) agree on every propositional variable occurring in \(\varphi\), then they assign the same truth value to \(\varphi\):
\[
v|_{\operatorname{Var}(\varphi)}
=
w|_{\operatorname{Var}(\varphi)}
\Longrightarrow
\widehat v(\varphi)=\widehat w(\varphi).
\]

\noindent\hyperref[prf:coincidence-lemma-propositional-logic]{Go to proof.}
\end{lemma}
\end{tcolorbox}

\begin{remark*}[Standard quantified statement]
\[
(\forall\varphi\in\WFF)
(\forall v,w:\Prop\to\mathbb{B})
\left[
(\forall P\in\operatorname{Var}(\varphi))(v(P)=w(P))
\to
\widehat v(\varphi)=\widehat w(\varphi)
\right].
\]
\end{remark*}

\begin{remark*}[Interpretation]
The truth value of a formula depends only on the variables actually occurring in that formula.
\end{remark*}

\begin{dependencies}
\begin{itemize}
  \item \hyperref[def:formula-variable-set-propositional-logic]{Variable Set of a Formula}
  \item \hyperref[thm:unique-extension-truth-assignment-propositional-logic]{Unique Extension of a Truth Assignment}
  \item \hyperref[thm:structural-induction-propositional-formulas]{Structural Induction for Propositional Formulas}
\end{itemize}
\end{dependencies}
% END GENERATED ARTIFACT: lem:coincidence-lemma-propositional-logic
% BEGIN GENERATED ARTIFACT: thm:finite-truth-table-propositional-logic
\begin{tcolorbox}[colback=thmbox, colframe=thmborder, arc=2pt,
  left=6pt, right=6pt, top=4pt, bottom=4pt,
  title={\small\textbf{Theorem (Finite Truth-Table Theorem)}},
  fonttitle=\small\bfseries]
\begin{theorem}[Finite Truth-Table Theorem]
\label{thm:finite-truth-table-propositional-logic}
Every propositional formula has a finite truth table. More precisely, if \(\operatorname{Var}(\varphi)\) has \(n\) elements, then \(\varphi\) has a truth table with \(2^n\) rows.

\noindent\hyperref[prf:finite-truth-table-propositional-logic]{Go to proof.}
\end{theorem}
\end{tcolorbox}

\begin{remark*}[Standard quantified statement]
\[
(\forall\varphi\in\WFF)
\left[
|\operatorname{Var}(\varphi)|=n
\to
\operatorname{TruthTable}(\varphi)\text{ has }2^n\text{ rows}
\right].
\]
\end{remark*}

\begin{remark*}[Interpretation]
Truth tables are finite because every formula contains only finitely many propositional variables.
\end{remark*}

\begin{dependencies}
\begin{itemize}
  \item \hyperref[def:truth-table-propositional-logic]{Truth Table}
  \item \hyperref[lem:finiteness-of-variables-propositional-formulas]{Finiteness of Variables in a Formula}
\end{itemize}
\end{dependencies}
% END GENERATED ARTIFACT: thm:finite-truth-table-propositional-logic
% BEGIN GENERATED ARTIFACT: def:boolean-function-propositional-logic
\begin{tcolorbox}[colback=propbox, colframe=propborder, arc=2pt,
  left=6pt, right=6pt, top=4pt, bottom=4pt,
  title={\small\textbf{Definition (Boolean Function)}},
  fonttitle=\small\bfseries]
\begin{definition}[Boolean Function]
\label{def:boolean-function-propositional-logic}
An \(n\)-ary Boolean function is a function
\[
f:\mathbb{B}^n\to\mathbb{B}.
\]
\end{definition}
\end{tcolorbox}

\begin{remark*}[Standard quantified statement]
\[
BooleanFunction_n(f)
\Longleftrightarrow
f:\mathbb{B}^n\to\mathbb{B}.
\]
\end{remark*}

\begin{remark*}[Definition predicate reading]
\[
BooleanFunction_n(f)
\]
means that \(f\) takes \(n\) truth values as input and returns one truth value.
\end{remark*}

\begin{remark*}[Interpretation]
Boolean functions are the abstract truth-functions that formulas may represent.
\end{remark*}

\begin{dependencies}
\begin{itemize}
  \item \hyperref[def:truth-value-propositional-logic]{Truth Value}
\end{itemize}
\end{dependencies}
% END GENERATED ARTIFACT: def:boolean-function-propositional-logic
% BEGIN GENERATED ARTIFACT: def:truth-function-induced-by-formula
\begin{tcolorbox}[colback=propbox, colframe=propborder, arc=2pt,
  left=6pt, right=6pt, top=4pt, bottom=4pt,
  title={\small\textbf{Definition (Truth Function Induced by a Formula)}},
  fonttitle=\small\bfseries]
\begin{definition}[Truth Function Induced by a Formula]
\label{def:truth-function-induced-by-formula}
Let \(\varphi\in\WFF\), and suppose \(\operatorname{Var}(\varphi)\subseteq\{P_1,\ldots,P_n\}\). The truth function induced by \(\varphi\) is the Boolean function
\[
f_\varphi:\mathbb{B}^n\to\mathbb{B}
\]
defined by
\[
f_\varphi(b_1,\ldots,b_n)=\widehat v(\varphi),
\]
where \(v(P_i)=b_i\) for each \(i\).
\end{definition}
\end{tcolorbox}

\begin{remark*}[Standard quantified statement]
\[
f_\varphi(b_1,\ldots,b_n)=\widehat v(\varphi)
\quad
\text{whenever }
v(P_i)=b_i\text{ for }1\le i\le n.
\]
\end{remark*}

\begin{remark*}[Definition predicate reading]
\[
\operatorname{TruthFunctionOfFormula}(f_\varphi,\varphi)
\]
means that \(f_\varphi\) records the truth behavior of \(\varphi\).
\end{remark*}

\begin{remark*}[Interpretation]
The induced truth function packages the truth table of a formula as a Boolean function.
\end{remark*}

\begin{dependencies}
\begin{itemize}
  \item \hyperref[def:boolean-function-propositional-logic]{Boolean Function}
  \item \hyperref[def:formula-variable-set-propositional-logic]{Variable Set of a Formula}
  \item \hyperref[thm:unique-extension-truth-assignment-propositional-logic]{Unique Extension of a Truth Assignment}
  \item \hyperref[lem:coincidence-lemma-propositional-logic]{Coincidence Lemma for Propositional Logic}
\end{itemize}
\end{dependencies}
% END GENERATED ARTIFACT: def:truth-function-induced-by-formula
% BEGIN GENERATED ARTIFACT: thm:counting-boolean-functions
\begin{tcolorbox}[colback=thmbox, colframe=thmborder, arc=2pt,
  left=6pt, right=6pt, top=4pt, bottom=4pt,
  title={\small\textbf{Theorem (Counting Boolean Functions)}},
  fonttitle=\small\bfseries]
\begin{theorem}[Counting Boolean Functions]
\label{thm:counting-boolean-functions}
There are exactly
\[
2^{2^n}
\]
Boolean functions \(\mathbb{B}^n\to\mathbb{B}\).

\noindent\hyperref[prf:counting-boolean-functions]{Go to proof.}
\end{theorem}
\end{tcolorbox}

\begin{remark*}[Standard quantified statement]
\[
\left|\{f:\mathbb{B}^n\to\mathbb{B}\}\right|=2^{2^n}.
\]
\end{remark*}

\begin{remark*}[Interpretation]
The domain \(\mathbb{B}^n\) has \(2^n\) elements. A Boolean function chooses one of two truth values for each input row.
\end{remark*}

\begin{dependencies}
\begin{itemize}
  \item \hyperref[def:boolean-function-propositional-logic]{Boolean Function}
\end{itemize}
\end{dependencies}
% END GENERATED ARTIFACT: thm:counting-boolean-functions
